\section{Axes d'amélioration}
\competence{C4.3.1}{Proposer des axes d'amélioration}{La présentation des recommandations argumentées d'amélioration}
\livrable{Des recommandations argumentées, évaluant les gains (coût, délai), réalistes et renforçant l'attractivité du logiciel.}

\subsection{Analyse des indicateurs et des retours utilisateurs}
\todo{Exploiter les indicateurs de performance (issus de la supervision, chap. 1) et les retours utilisateurs (support, recette) pour identifier les points d'amélioration.}

\subsection{Recommandations argumentées}
\todo{Présenter 3 à 4 recommandations ; pour chacune : problème adressé, gain attendu, coût et délai de mise en œuvre, faisabilité au regard du projet.}

\subsection{Priorisation et attractivité}
\todo{Prioriser (matrice valeur/effort) et relier chaque recommandation au renforcement de l'attractivité du logiciel.}
